% Editable vector source; requires tikz with arrows.meta and positioning.
\begin{figure*}[t]
\centering
\begin{tikzpicture}[
  font=\sffamily\small,
  box/.style={draw=black!45,rounded corners=2pt,align=center,inner sep=7pt},
  flow/.style={-{Latex[length=2mm]},thick,draw=black!65},
  note/.style={align=center,font=\sffamily\footnotesize},
  heading/.style={anchor=west,font=\sffamily\bfseries\small}
]
\node[heading] at (0,0) {(a) Build the phenotype-query graph};
\node[box,fill=blue!5,text width=4.9cm] (graph) at (2.6,-1.25)
  {Candidate genes + phenotype hierarchy\\[2pt]
   Gene--phenotype and gene--function edges\\
   Typed links from phenotype definitions};
\node[note,text width=4.9cm] at (2.6,-2.55)
  {Training: gene--phenotype links\\
   Query: disease phenotype set};
\node[box,fill=blue!5,text width=4.9cm] (encoder) at (2.6,-3.8)
  {Pretrained relational graph network\\[2pt]
   Evaluate initial and adapted models;\\select on validation diseases};
\draw[flow] (graph.west) -- ++(-.45,0) |- (encoder.west);

\node[heading] at (6.1,0) {(b) Rank genes for a disease phenotype set};
\node[box,text width=4.0cm] (query) at (8.5,-1.0)
  {Unseen disease\\$P_d=\{p_1,\ldots,p_m\}$};
\node[box,fill=blue!5,text width=4.0cm] (atomic) at (8.5,-2.65)
  {Project each phenotype to all genes\\$s(g,p_1),\ldots,s(g,p_m)$};
\node[box,fill=green!6,text width=4.0cm] (conjunction) at (8.5,-4.35)
  {Combine evidence by conjunction\\[2pt]
   $q(g,d)=\frac{1}{m}\sum_{p\in P_d}\log\sigma(s(g,p))$};
\draw[flow] (query) -- (atomic);
\draw[flow] (encoder.east) -- ++(.5,0) |- (atomic.west);
\draw[flow] (atomic) -- (conjunction);

\node[box,fill=green!6,text width=3.0cm] (calibration) at (13.1,-2.65)
  {Compare with the gene's\\reference-disease scores\\[2pt]
   $z(g,d)=\dfrac{q(g,d)-\mu_g}{\sigma_g+\epsilon}$};
\node[box,text width=3.0cm] (ranking) at (13.1,-4.35)
  {Rank all candidate genes\\[2pt]
   Retain raw and calibrated results};
\draw[flow] (conjunction.east) -- ++(.45,0) |- (calibration.west);
\draw[flow] (calibration) -- (ranking);

\node[heading] at (0,-5.75) {(c) Audit the scoring rule and evaluation population};
\node[box,text width=4.25cm] at (2.4,-6.75)
  {Training-association status\\[2pt]
   Seen/unseen strata and within-status ranks};
\node[box,text width=4.25cm] at (7.55,-6.75)
  {Training-phenotype overlap\\[2pt]
   Shared phenotypes and zero-overlap cohort};
\node[box,text width=4.25cm] at (12.7,-6.75)
  {Repeated profiles and positives\\[2pt]
   Reference deduplication, co-causal filtering and profile weighting};
\end{tikzpicture}
\caption{PhenoQuery and its evaluation.
(a) The typed PhenoQuery graph combines gene annotations, ontology hierarchy
and typed Entity--Quality links extracted from phenotype definitions. Its
standard construction omits training disease nodes and expression edges;
a separate disease-bridge arm tests the effect of retaining training diseases.
(b) Each query phenotype supplies an atomic gene-link score. Length-normalized
log-sigmoid conjunction combines these scores, and per-gene reference
calibration adjusts their baseline level and scale. The query disease is
excluded from its calibration reference panel.
(c) Training status, phenotype overlap and repeated profiles define distinct
score adjustments or evaluation sensitivities. Their effects are reported
separately, with each analysis retaining its defined scoring rule and population.}
\label{fig:phenoquery}
\end{figure*}
